\newif\ifglobal

%%%%%%%%%%%%%%%%%%%%%%%
%%% Compile to view %%%
%%%%%%%%%%%%%%%%%%%%%%%

%%%%%%%%%%%%%%%%%%%%%%%%%%%%%%%%%%%%%%%%%%%%%%%%%%%
%%% Readme:                                     %%%
%%% https://www.overleaf.com/read/bwdjvfjksvjy/ %%%
%%%%%%%%%%%%%%%%%%%%%%%%%%%%%%%%%%%%%%%%%%%%%%%%%%%

% >>> M?: marks a module setting number or important notice

% >>> M4: Set {./relative-path} to external files for this module <<<
\def \modulefiles {./13-PID}
% To add an external file, use:
% (Replace '---' with actual filename)
% '\input{\modulefiles/tables/---.tex}' for tables
% '\includegraphics{\modulefiles/figures/---}' for figures
% '\subfile{\modulefiles/---}' for register files

% >>> M1: This module is in global project? <<<
% 'YES' - uncomment; 'NO' - comment out
\globaltrue

\ifglobal
    \documentclass[main\_\_EN.tex]{subfiles}

\else
    \documentclass[openany, 10pt]{book}

    % >>> M2: In ./00-shared/preamble-trm-module, also find
    %         settings for visibility of todo notes, watermark <<<
    \usepackage{preamble-trm-module}

    % >>> M3: Temporary labels in English,
    %         you can add your own labels there <<<
    \myexternaldocument{./00-shared/temp-labels-en}

\fi %\ifglobal


\setcounter{secnumdepth}{4}


% Define formatting of parts and chapters in text
\setenglishformatting


% >>> M5: If this module needs any updates in
% './00-shared/preamble-trm-module.sty'
% (include additional package, etc.), contact your module PM
% or create the file './00-shared/changelog.tex' make the
% necessary updates yourself and document them in this file <<<

% Make text left-aligned and allow latex to break pages naturally, comment out for full justification
\raggedright
\raggedbottom


\begin{document}

\hypertarget{pidctrl}{}
\chapter{Process ID Controller (PID)}
\label{mod:pid}

\section{Overview}

The ESP32 is a dual core device and is capable of running and managing multiple processes. The PID Controller supports switching of PID when a process switch occurs. In addition to PID management, the PID Controller also facilitates management of nested interrupts by recording execution status just before an interrupt service routine is executed. This enables the user application to manage process switches and nested interrupts more efficiently.

\section{Features}

The PID Controller features:
\begin{itemize}
    \item Process management and priority
    \item Process PID switch
    \item Interrupt information recording
    \item Nested interrupt management
\end{itemize}


\section{Functional Description}

Eight processes run on the CPU, and are assigned with PID of 0 \verb+~+ 7 respectively.
Among the eight processes, processes with PID of 0 or 1 are elevated processes with higher authority compared to processes with PID ranging from 2 \verb+~+ 7.

A CPU process switch may occur in two cases:
\begin{itemize}
    \item An interrupt occurs and the CPU fetches an instruction from the interrupt vector. Instruction fetch or execution from interrupt vector is always treated as a process with PID of 0, irrespective of which process was being executed on the CPU when the interrupt occurred.
    \item A currently active process explicitly performs a process switch. Only elevated processes with PID of 0 or 1 may perform a process switch.
\end{itemize}
\subsection{Interrupt Identification}

Interrupts are classified into seven priority levels:
Level 1, Level 2, Level 3, Level 4, Level 5, Level 6 (Debug), and NMI.
Each level of interrupt is assigned an interrupt vector entry address.
The PID Controller recognizes CPU instruction fetch from an interrupt vector entry address and automatically switches PID to 0. If CPU only accesses the interrupt vector entry address, PID Controller performs no action.

\hyperref[regdesc:PIDCTRLINTERRUPTENABLEREG]{PIDCTRL\_INTERRUPT\_ENABLE\_REG} determines whether the PID Controller identifies and registers an interrupt of certain priority. When a bit of register \hyperref[regdesc:PIDCTRLINTERRUPTENABLEREG]{PIDCTRL\_INTERRUPT\_ENABLE\_REG} is 1, PID Controller will take action when CPU fetches instruction from the interrupt vector entry address of the corresponding interrupt. Otherwise, PID Controller performs no action.
The registers PIDCTRL\_INTERRUPT\_ADDR\_1\_REG \verb+~+ PIDCTRL\_INTERRUPT\_ADDR\_7\_REG define the interrupt vector entry address for all the interrupt priority levels.
For details please refer to Table \ref{table:Interrupt}.


\begin{table}[H]
    \centering
    \caption{Interrupt Vector Entry Address}
    \label{table:Interrupt}
    \begin{tabular}{|p{2.5cm}|p{6.5cm}|p{6cm}|}
        \hline
        \rowcolor{lightgray}
        Priority level &
        \makecell[l]{PIDCTRL\_INTERRUPT\_ENABLE\_REG bit\\controlling interrupt identification} &
        Interrupt vector entry address \\
        \hline
        Level 1 & 1 & \hyperref[regdesc:PIDCTRLINTERRUPTADDR1REG]{PIDCTRL\_INTERRUPT\_ADDR\_1\_REG}\\
        \hline
        Level 2 & 2 & \hyperref[regdesc:PIDCTRLINTERRUPTADDR2REG]{PIDCTRL\_INTERRUPT\_ADDR\_2\_REG} \\
        \hline
        Level 3 & 3 & \hyperref[regdesc:PIDCTRLINTERRUPTADDR3REG]{PIDCTRL\_INTERRUPT\_ADDR\_3\_REG}\\
        \hline
        Level 4 & 4 & \hyperref[regdesc:PIDCTRLINTERRUPTADDR4REG]{PIDCTRL\_INTERRUPT\_ADDR\_4\_REG}\\
        \hline
        Level 5 & 5 & \hyperref[regdesc:PIDCTRLINTERRUPTADDR5REG]{PIDCTRL\_INTERRUPT\_ADDR\_5\_REG}\\
        \hline
        Level 6 ( Debug ) & 6 & \hyperref[regdesc:PIDCTRLINTERRUPTADDR6REG]{PIDCTRL\_INTERRUPT\_ADDR\_6\_REG}\\
        \hline
        NMI & 7 & \hyperref[regdesc:PIDCTRLINTERRUPTADDR7REG]{PIDCTRL\_INTERRUPT\_ADDR\_7\_REG}\\
        \hline
    \end{tabular}
\end{table}


\subsection{Information Recording}

When PID Controller identifies an interrupt, it records three items of information in addition to switching PID to 0. The recorded information includes the priority level of current interrupt, previous interrupt status of the system and the previous process running on the CPU.

PID Controller records the priority level of the current interrupt in register \hyperref[regdesc:PIDCTRLLEVELREG]{PIDCTRL\_LEVEL\_REG}. For details please refer to Table \ref{table:PIDCTRL_LEVEL_REG}.

\begin{table}[H]
    \centering
    \caption{Configuration of PIDCTRL\_LEVEL\_REG}
    \label{table:PIDCTRL_LEVEL_REG}
    \begin{tabular}{|p{5cm}|p{10cm}|}
        \hline
        \rowcolor{lightgray}
        Value& Priority level of the current interrupt \\
        \hline
        0 & No interrupt \\
        \hline
        1 & Level 1  \\
        \hline
        2 & Level 2  \\
        \hline
        3 & Level 3  \\
        \hline
        4 & Level 4  \\
        \hline
        5 & Level 5  \\
        \hline
        6 & Level 6  \\
        \hline
        7 & NMI \\
        \hline
    \end{tabular}
\end{table}

PID Controller also records in register \hyperref[regdesc:PIDCTRLFROMnREG]{PIDCTRL\_FROM\_\textcolor{red}{\it{n}}\_REG} the status of the system before the interrupt occurred.
The bit width of register \hyperref[regdesc:PIDCTRLFROMnREG]{PIDCTRL\_FROM\_\textcolor{red}{\it{n}}\_REG} is 7. The highest four bits represent the interrupt status of the system before the interrupt indicated by the register occurred. The lowest three bits represent the process running on the CPU before the interrupt indicated by the register occurred.
For details please refer to Table \ref{table:PIDCTRL_FROM_n_REG}.

\begin{table}[H]
    \centering
    \caption{Configuration of PIDCTRL\_FROM\_\textcolor{red}{\it{n}}\_REG}
    \label{table:PIDCTRL_FROM_n_REG}
    \begin{tabular}{|p{2cm}|p{5cm}|}
        \hline
        \rowcolor{lightgray}
        \verb+[+6:3\verb+]+ & Previous interrupt \\
        \hline
        0 & No interrupt \\
        \hline
        1 & Level 1 Interrupt \\
        \hline
        2 & Level 2 Interrupt \\
        \hline
        3 & Level 3 Interrupt \\
        \hline
        4 & Level 4 Interrupt \\
        \hline
        5 & Level 5 Interrupt \\
        \hline
        6 & Level 6 Interrupt \\
        \hline
        7 & Level 7 Interrupt \\
        \hline
    \end{tabular}
    \begin{tabular}{|p{2cm}|p{5cm}|}
        \hline
        \rowcolor{lightgray}
        \verb+[+2:0\verb+]+ & Previous process \\
        \hline
        0 & Process with PID of 0 \\
        \hline
        1 & Process with PID of 1 \\
        \hline
        2 & Process with PID of 2 \\
        \hline
        3 & Process with PID of 3 \\
        \hline
        4 & Process with PID of 4 \\
        \hline
        5 & Process with PID of 5 \\
        \hline
        6 & Process with PID of 6 \\
        \hline
        7 & Process with PID of 7 \\
        \hline
    \end{tabular}
\end{table}


PID Controller possesses registers PIDCTRL\_FROM\_1\_REG \verb+~+ PIDCTRL\_FROM\_7\_REG, which correspond to the interrupts of Level 1, Level 2, Level 3, Level 4, Level 5, Level 6 (Debug), and NMI respectively. This enables the system to implement interrupt nesting. Please refer to Table \ref{figure:Interrupt_Nesting} for examples.

\begin{figure}[!ht]
    \centering
    \includegraphics[width=1.0\textwidth]{\modulefiles/figures/Interrupt-Nesting_en.png}
    \caption{Interrupt Nesting}
    \label{figure:Interrupt_Nesting}
\end{figure}

If the configuration of register \hyperref[regdesc:PIDCTRLINTERRUPTENABLEREG]{PIDCTRL\_INTERRUPT\_ENABLE\_REG} prevents PID Controller from identifying an interrupt, PID Controller will not record any information, and \hyperref[regdesc:PIDCTRLLEVELREG]{PIDCTRL\_LEVEL\_REG} and \hyperref[regdesc:PIDCTRLFROMnREG]{PIDCTRL\_FROM\_\textcolor{red}{\it{n}}\_REG} will remain unchanged.

\subsection{Proactive Process Switching}

As mentioned before, only an elevated process with PID of 0/1 can initiate a process switch.
The new process may have any PID from 0 \verb+~+ 7 after the process switch.
The key for successful proactive process switching is that when the last command of the current process switches to the first command of the new process, PID should switch from 0/1 to that of the new process.

The software procedure for proactive process switching is as follows:

\newcounter{STEP_MASK_NON_NMI}
\newcounter{STEP_MASK_NMI}
\newcounter{STEP_SET_DELAY}
\newcounter{STEP_SET_PID_NEW}
\newcounter{STEP_RESTORE}
\newcounter{STEP_CHANGE}
\newcounter{STEP_RELEASE_NON_NMI}
\newcounter{STEP_JUMP}
\begin{enumerate}
    \item Mask all the interrupts except NMI by using software.
        \setcounter{STEP_MASK_NON_NMI}{\theenumi}
    \item Set register \hyperref[regdesc:PIDCTRLNMIMASKENABLEREG]{PIDCTRL\_NMI\_MASK\_ENABLE\_REG} to 1 to generate a CPU NMI Interrupt Mask signal.
        \setcounter{STEP_MASK_NMI}{\theenumi}
    \item Configure registers \hyperref[regdesc:PIDCTRLPIDDELAYREG]{PIDCTRL\_PID\_DELAY\_REG} and \hyperref[regdesc:PIDCTRLNMIDELAYREG]{PIDCTRL\_NMI\_DELAY\_REG}.
        \setcounter{STEP_SET_DELAY}{\theenumi}
    \item Configure register \hyperref[regdesc:PIDCTRLPIDNEWREG]{PIDCTRL\_PID\_NEW\_REG}.
        \setcounter{STEP_SET_PID_NEW}{\theenumi}
    \item Configure register \hyperref[regdesc:PIDCTRLLEVELREG]{PIDCTRL\_LEVEL\_REG} and \hyperref[regdesc:PIDCTRLFROMnREG]{PIDCTRL\_FROM\_\textcolor{red}{\it{n}}\_REG}.
        \setcounter{STEP_RESTORE}{\theenumi}
    \item Set register \hyperref[regdesc:PIDCTRLPIDCONFIRMREG]{PIDCTRL\_PID\_CONFIRM\_REG} and register \hyperref[regdesc:PIDCTRLNMIMASKDISABLEREG]{PIDCTRL\_NMI\_MASK\_DISABLE\_REG} to 1.
        \setcounter{STEP_CHANGE}{\theenumi}
    \item Revoke the masking of all interrupts but NMI.
        \setcounter{STEP_RELEASE_NON_NMI}{\theenumi}
    \item Switch to the new process and fetch instruction.
        \setcounter{STEP_JUMP}{\theenumi}
\end{enumerate}

Though we can deal with interrupt nesting,
an elevated process should not be interrupted during the process switching, and therefore the interrupts have been masked in step \arabic{STEP_MASK_NON_NMI} and step \arabic{STEP_MASK_NMI}.

In step \arabic{STEP_SET_DELAY}, the configured values of registers \hyperref[regdesc:PIDCTRLPIDDELAYREG]{PIDCTRL\_PID\_DELAY\_REG} and \hyperref[regdesc:PIDCTRLNMIDELAYREG]{PIDCTRL\_NMI\_DELAY\_REG} will affect step \arabic{STEP_CHANGE}.

In step \arabic{STEP_SET_PID_NEW}, the configured value of register
\hyperref[regdesc:PIDCTRLPIDNEWREG]{PIDCTRL\_PID\_NEW\_REG} will be the new PID after step \arabic{STEP_CHANGE}.

If the system is currently in a nested interrupt and needs to revert to the previous interrupt, register \hyperref[regdesc:PIDCTRLLEVELREG]{PIDCTRL\_LEVEL\_REG} must be restored based on the information recorded in register \hyperref[regdesc:PIDCTRLFROMnREG]{PIDCTRL\_FROM\_\textcolor{red}{\it{n}}\_REG} in step \arabic{STEP_RESTORE}.

In step \arabic{STEP_CHANGE}, after the values of register \hyperref[regdesc:PIDCTRLPIDCONFIRMREG]{PIDCTRL\_PID\_CONFIRM\_REG} and register \hyperref[regdesc:PIDCTRLNMIMASKDISABLEREG]{PIDCTRL\_NMI\_MASK\_DISABLE\_REG} are set to 1, PID Controller will not immediately switch PID to the value of register \hyperref[regdesc:PIDCTRLPIDNEWREG]{PIDCTRL\_PID\_NEW\_REG}, nor disable CPU NMI Interrupt Mask signal at once. Instead, PID Controller performs each task after a different number of clock cycles. The numbers of clock cycles are the values specified in register \hyperref[regdesc:PIDCTRLPIDDELAYREG]{PIDCTRL\_PID\_DELAY\_REG} and \hyperref[regdesc:PIDCTRLNMIDELAYREG]{PIDCTRL\_NMI\_DELAY\_REG} respectively.

In step \arabic{STEP_RELEASE_NON_NMI}, other tasks can be implemented as well.
To do this, the cost of those tasks should be included when configuring registers \hyperref[regdesc:PIDCTRLPIDDELAYREG]{PIDCTRL\_PID\_DELAY\_REG} and \hyperref[regdesc:PIDCTRLNMIDELAYREG]{PIDCTRL\_NMI\_DELAY\_REG} in step \arabic{STEP_SET_DELAY}.

\hypertarget{pidctrl-reg-summ}{}
\section{Register Summary}\label{sec:pidctrl-reg-summ}

\subfile{\modulefiles/pidc_reg-summary_en}

\section{Registers}

The addresses in parenthesis besides register names are the register addresses relative to the PID Controller base address provided in Table \ref{table:Peripheral_Address_Mapping} \textit{Peripheral Address Mapping} in Chapter \ref{mod:sysmem} \textit{\nameref{mod:sysmem}}. The absolute register addresses are listed in Section \ref{sec:pidctrl-reg-summ} \textit{\nameref{sec:pidctrl-reg-summ}}. 

\subfile{\modulefiles/pidc_reg_en}

\end{document}
